% Options for packages loaded elsewhere
\PassOptionsToPackage{unicode}{hyperref}
\PassOptionsToPackage{hyphens}{url}
%
\documentclass[
  11pt,
  letterpaper,
]{article}
\usepackage{amsmath,amssymb}
\usepackage{iftex}
\ifPDFTeX
  \usepackage[T1]{fontenc}
  \usepackage[utf8]{inputenc}
  \usepackage{textcomp} % provide euro and other symbols
\else % if luatex or xetex
  \usepackage{unicode-math} % this also loads fontspec
  \defaultfontfeatures{Scale=MatchLowercase}
  \defaultfontfeatures[\rmfamily]{Ligatures=TeX,Scale=1}
\fi
\usepackage{lmodern}
\pdfmapfile{+lm.map}
\pdfmapfile{+symbols.map}
\pdfmapfile{+cm.map}
\ifPDFTeX\else
  % xetex/luatex font selection
\fi
% Use upquote if available, for straight quotes in verbatim environments
\IfFileExists{upquote.sty}{\usepackage{upquote}}{}
\IfFileExists{microtype.sty}{% use microtype if available
  \usepackage[]{microtype}
  \UseMicrotypeSet[protrusion]{basicmath} % disable protrusion for tt fonts
}{}
\makeatletter
\@ifundefined{KOMAClassName}{% if non-KOMA class
  \IfFileExists{parskip.sty}{%
    \usepackage{parskip}
  }{% else
    \setlength{\parindent}{0pt}
    \setlength{\parskip}{6pt plus 2pt minus 1pt}}
}{% if KOMA class
  \KOMAoptions{parskip=half}}
\makeatother
\usepackage{xcolor}
\usepackage[margin=0.85in]{geometry}
\usepackage{longtable,booktabs,array}
\usepackage{calc} % for calculating minipage widths
% Correct order of tables after \paragraph or \subparagraph
\usepackage{etoolbox}
\makeatletter
\patchcmd\longtable{\par}{\if@noskipsec\mbox{}\fi\par}{}{}
\makeatother
% Allow footnotes in longtable head/foot
\IfFileExists{footnotehyper.sty}{\usepackage{footnotehyper}}{\usepackage{footnote}}
\makesavenoteenv{longtable}
\setlength{\emergencystretch}{3em} % prevent overfull lines
\providecommand{\tightlist}{%
  \setlength{\itemsep}{0pt}\setlength{\parskip}{0pt}}
\setcounter{secnumdepth}{-\maxdimen} % remove section numbering
\ifLuaTeX
\usepackage[bidi=basic]{babel}
\else
\usepackage[bidi=default]{babel}
\fi
\babelprovide[main,import]{american}
% get rid of language-specific shorthands (see #6817):
\let\LanguageShortHands\languageshorthands
\def\languageshorthands#1{}
\ifLuaTeX
  \usepackage{selnolig}  % disable illegal ligatures
\fi
\usepackage{bookmark}
\IfFileExists{xurl.sty}{\usepackage{xurl}}{} % add URL line breaks if available
\urlstyle{same}
\hypersetup{
  pdftitle={Fourier endpoint minimax and autocorrelation deficit rank bounds},
  pdflang={en-US},
  hidelinks,
  pdfcreator={LaTeX via pandoc}}

\title{Fourier endpoint minimax and autocorrelation deficit rank bounds}
\usepackage{etoolbox}
\makeatletter
\providecommand{\subtitle}[1]{% add subtitle to \maketitle
  \apptocmd{\@title}{\par {\large #1 \par}}{}{}
}
\makeatother
\subtitle{A two theorem working manuscript}
\author{}
\date{11 October 2026}

\begin{document}
\maketitle

\section{Abstract}\label{abstract}

This manuscript gives complete proofs of two previously established
project results, keeping their hypotheses and conclusions separate. For
the cyclic Fourier matrix of order \(N=2^m\geq8\), the minimum, over
independent row and column permutations, of the maximum strict
off-diagonal corner rank is \(N/4+1\). Double bit reversal attains the
minimum. The minimum rank of an arbitrary half-size square is instead
\(N/4\). Separately, for a unitary character table of any finite abelian
group, a complementary Gram-matrix second moment gives a rectangular
rank bound and a positive low-rank approximation gap in terms of the
maximum nonzero-translation overlap of the column set. All degenerate
sizes, normalization choices, and singular-value indices are specified
below.

This is a consolidation of accepted proof sources, not a new theorem or
fusion. Digit reversal, Fourier support uncertainty, the flat-unitary
rank bound, translation-overlap statistics, and spectral approximation
arguments have established antecedents. The exact scope's historical
priority remains unresolved. No general FFT arithmetic-complexity
conclusion or publication-acceptance claim is made.

\section{1 Scope and conventions}\label{scope-and-conventions}

For a finite set \(E\), write \(\mathbb C^E\) for its coordinate space
with the standard Hermitian inner product. If \(A\) has row set \(E\)
and column set \(D\), then \(A[R,H]\) is the submatrix with
\(R\subseteq E\) and \(H\subseteq D\). Reordering the chosen rows or
columns does not change rank or singular values. Complements are taken
in the full relevant indexing set. An empty matrix has rank zero.

For \(N\geq1\), put \([N]=\{0,\ldots,N-1\}\). If \(\rho\) is a
permutation of \([N]\), use the convention \((P_\rho x)_i=x_{\rho(i)}\).
Thus

\[
(P_\rho A P_\sigma^{T})_{ij}=A_{\rho(i),\sigma(j)}.
\]

For an \(N\times N\) matrix define

\[
q(A)=\max_{1\leq k<N}\max\bigl\{
\operatorname{rank}A[0:k,k:N],\,
\operatorname{rank}A[k:N,0:k]\bigr\}.                 \tag{1}
\]

The interval notation is half-open. These two blocks are the strict
off-diagonal corners across the same coordinate cut. Set \(q(A)=0\) for
\(N=1\), when there is no cut. The two endpoint permutations in the
minimization below are arbitrary and independent.

The unnormalized cyclic Fourier matrix and its unitary normalization are

\[
F_N=(\omega_N^{ab})_{a,b\in[N]},\qquad
\omega_N=e^{2\pi i/N},\qquad U_N=N^{-1/2}F_N.        \tag{2}
\]

For a finite abelian group \(G\) of order \(N\), written additively, its
dual group is \(\widehat G\). The unitary character table used here is

\[
U_{\chi,x}=N^{-1/2}\chi(x),\qquad
\chi\in\widehat G,\ x\in G.                       \tag{3}
\]

Characters have modulus one. Character orthogonality makes \(U\)
unitary. Conjugating the Fourier convention changes none of the
conclusions. Nonzero scalar normalization does not change exact rank; it
does change singular values and operator-norm errors. Throughout the
analytic statements, \(\|\cdot\|_2\) denotes the operator norm for
Euclidean norms, rather than the Frobenius norm.

\subsection{Exact assumption and convention
table}\label{exact-assumption-and-convention-table}

\begin{longtable}[]{@{}
  >{\raggedright\arraybackslash}p{(\columnwidth - 2\tabcolsep) * \real{0.2600}}
  >{\raggedright\arraybackslash}p{(\columnwidth - 2\tabcolsep) * \real{0.7400}}@{}}
\toprule\noalign{}
\begin{minipage}[b]{\linewidth}\raggedright
Item
\end{minipage} & \begin{minipage}[b]{\linewidth}\raggedright
Convention and required hypothesis
\end{minipage} \\
\midrule\noalign{}
\endhead
\bottomrule\noalign{}
\endlastfoot
Endpoint minimax & Cyclic \(N=2^m\); independent endpoint permutations;
one matrix and a common cut coordinate order \\
Exact field scope in Theorem 1 & Complex numbers, or a
characteristic-zero field containing a primitive \(N\)th root; no square
root of \(N\) is needed for the unnormalized statement \\
Deficit theorem & Any finite abelian \(G\), \(N\geq2\); complex
unitary-normalized character table; \(1\leq r,h\leq N\) \\
Selected sets & Rows \(R\subseteq\widehat G\), columns \(H\subseteq G\);
\(r=|R|\), \(h=|H|\) \\
Translation overlap & \(c_H(s)=|H\cap(H-s)|=|\{x\in H:x+s\in H\}|\) \\
Deficit & \(d_H=h-\max_{s\neq0}c_H(s)\); defined only when \(N\geq2\) \\
Singular values & Descending square roots of the \(h\) eigenvalues of
\(V^*V\), including zeros if \(h>r\) \\
Approximation rank & \(W\in\mathbb C^{r\times h}\), integer \(t\geq0\),
\(\operatorname{rank}W\leq t\); strict gap condition \(t<h-B_H\) \\
Half-square result & Cyclic \(N=4L\); arbitrary row and column sets,
each of size \(N/2\); \(N=2\) treated separately \\
Finite characteristic & No general lower-bound or positive-gap transfer
is asserted; modular checks are controls only \\
Open extensions & Noncyclic matching minimax upper bounds and exact
all-size minimum rank sums are not assumptions of either theorem \\
\end{longtable}

\section{2 The cyclic dyadic endpoint
theorem}\label{the-cyclic-dyadic-endpoint-theorem}

\textbf{Theorem 1.} Let \(N=2^m\), \(m\geq0\). Over \(\mathbb C\),

\[
\min_{\rho,\sigma\in S_N}q(P_\rho F_NP_\sigma^T)
=\begin{cases}
0,&N=1,\\
1,&N=2\text{ or }4,\\
N/4+1,&N\geq8.
\end{cases}                                                    \tag{4}
\]

For every size, bit reversal at both endpoints is a minimizing pair. The
same exact statement holds over every characteristic-zero field
containing a primitive \(N\)th root, with that root in place of
\(\omega_N\).

Separately, for \(N=2^m\geq4\),

\[
\min_{|R|=|H|=N/2}\operatorname{rank}F_N[R,H]=N/4.             \tag{5}
\]

For \(N=2\), this half-square minimum is \(1\). At \(N=1\) there is no
integer half-square size. Equation (5) is not a \(+1\) statement.

\subsection{2.1 The universal critical rectangle lower
bound}\label{the-universal-critical-rectangle-lower-bound}

Write \(N=4L\) with \(L\geq2\). Take arbitrary
\(R,H\subseteq\mathbb Z_N\) with

\[
|R|=2L-1,\qquad |H|=2L+1=:h.
\]

Let \(I=\mathbb Z_N\setminus R\), so \(|I|=h\), and let

\[
V=U_N[R,H],\qquad A=U_N[I,H]^*U_N[I,H].
\]

The columns of \(U_N\) are orthonormal, hence

\[
0\leq A\leq I_h,\qquad V^*V=I_h-A.                          \tag{6}
\]

Consequently \(\ker V\) is the eigenvalue-one eigenspace of \(A\). Its
dimension \(d\) is at most \(\operatorname{tr}(A^2)\), since every
eigenvalue of \(A\) belongs to \([0,1]\).

Set \(f_I(s)=\sum_{a\in I}\omega_N^{as}\). The \((x,y)\) entry of \(A\),
for \(x,y\in H\), is \(N^{-1}f_I(y-x)\). Grouping ordered pairs by their
difference gives

\[
\operatorname{tr}(A^2)=\frac1{N^2}\sum_{s\in\mathbb Z_N}
c_H(s)|f_I(s)|^2.                                             \tag{7}
\]

Here \(c_H(0)=h\). Every nonzero translation of \(\mathbb Z_{2^m}\) has
even order and partitions the group into cycles of that even length. An
invariant set is a union of those cycles and therefore has even
cardinality. Since \(h\) is odd, \(H\) is not invariant under any such
translation. Thus \(c_H(s)\leq h-1\) for \(s\neq0\). Orthogonality also
gives

\[
\sum_s|f_I(s)|^2=Nh,\qquad f_I(0)=h.
\]

Substitution into (7) yields

\[
\begin{aligned}
\operatorname{tr}(A^2)
&\leq\frac{h h^2+(h-1)(Nh-h^2)}{N^2}\\
&=L+\frac34+\frac1{4L}+\frac1{16L^2}=:B_L<L+1.              \end{aligned} \tag{8}
\]

For the last inequality, \(4L^2-4L-1>0\) for every integer \(L\geq2\).
Since \(d\) is an integer, \(d\leq L\). Therefore every such rectangle
satisfies

\[
\operatorname{rank}F_N[R,H]=h-d\geq L+1.                    \tag{9}
\]

After any independent row and column reorderings, the upper-right corner
at \(k=2L-1\) has exactly these dimensions. Equation (9) therefore
proves the lower bound in (4) uniformly over both permutations. No
assumption about the positions of either selected set was used.

\subsection{2.2 The exact bit reversal
recurrence}\label{the-exact-bit-reversal-recurrence}

Let \(b_m(j)\) be reversal of the \(m\) binary digits of \(j\), with
\(b_0(0)=0\). Define

\[
Q_N[i,j]=\omega_N^{b_m(i)b_m(j)},\qquad
r_N(k)=\operatorname{rank}Q_N[0:k,k:N].
\]

Put \(r_N(0)=r_N(N)=0\). We prove, for every dyadic \(N=4L\geq4\) and
\(0\leq k\leq2L\),

\[
r_N(k)=\lceil k/2\rceil+r_L(\lfloor k/2\rfloor).             \tag{10}
\]

The smaller table uses \(\omega_L=\omega_N^4\), including the one-entry
table when \(L=1\). The first \(2L\) bit-reversed labels are all the
even labels. Write them as \(2q_i\). If \(z_j=b_{m-2}(j)\), then

\[
q_{2j}=z_j,\qquad q_{2j+1}=L+z_j,\qquad 0\leq j<L.          \tag{11}
\]

Let \(p=\lfloor k/2\rfloor\) and \(d=\lceil k/2\rceil\). The selected
rows contain \(p\) complete pairs from (11), and one unpaired row if
\(k\) is odd. They have \(d\) distinct residues modulo \(L\).

All odd natural column labels \(2\ell+1\), \(0\leq\ell<2L\), lie to the
right of the cut. In row label \(2q\) their entries are

\[
\omega_N^{(2q)(2\ell+1)}
=\omega_{2L}^{q}\omega_L^{q\ell}.                           \tag{12}
\]

Rows \(q\) and \(q+L\) have opposite values on these columns. On one
representative of each distinct residue, removal of the nonzero phase
\(\omega_{2L}^q\) leaves \(d\) distinct rows of the full Fourier matrix
\(F_L\), with its columns repeated twice. These rows are independent, so
the odd-column span has dimension exactly \(d\).

Define \(\pi:\mathbb C^k\to\mathbb C^p\) by adding the two coordinates
in each complete pair and ignoring the unpaired coordinate, if present.
The odd-column span lies in \(\ker\pi\). Since \(\pi\) is surjective and
\(k-p=d\), both spaces have dimension \(d\) and hence are equal. Passing
to the quotient by the odd-column span is therefore precisely passing
through \(\pi\).

A surviving even column label \(2\ell\) has equal entries on each paired
row, and its image under \(\pi\) is

\[
2\bigl(\omega_L^{z_j\ell}\bigr)_{0\leq j<p}.                \tag{13}
\]

The column positions \(k,\ldots,2L-1\) in (11) give, modulo \(L\),
exactly the residues \(z_p,\ldots,z_{L-1}\), with possible repetitions.
At an odd cut \(k=2p+1\), the second member of the partially removed
pair still survives and supplies \(z_p\). Dropping it would change the
recurrence incorrectly. The quotient columns therefore have rank
\(r_L(p)\). Adding back the \(d\)-dimensional odd-column span proves
(10), including \(p=0\) and \(p=L\).

\subsection{2.3 Reflection and the maximum cut
rank}\label{reflection-and-the-maximum-cut-rank}

Digit reversal satisfies \(b_m(N-1-i)=N-1-b_m(i)\). Consequently

\[
Q_N[N-1-i,N-1-j]
=\omega_N^{1+b_m(i)+b_m(j)}Q_N[i,j].                       \tag{14}
\]

The factors on the right are invertible row and column scalings.
Simultaneous reversal carries an upper corner at \(k\) to a lower corner
at \(N-k\); \(Q_N\) is symmetric, so

\[
r_N(k)=r_N(N-k).                                           \tag{15}
\]

For \(k=2a\leq2L\), (10) and the column-count bound \(r_L(a)\leq L-a\)
give \(r_N(k)\leq L\). For \(k=2a+1<2L\), they give \(r_N(k)\leq L+1\).
Reflection covers every other cut. If \(L\geq2\), then \(r_L(L-1)=1\),
because this is a nonempty single-column Fourier block with nonzero
entries. Thus \(r_N(2L-1)=L+1\). Symmetry equates the upper and lower
ranks at each cut. This proves the matching upper bound for \(N\geq8\).

The bases are \(r_1(0)=r_1(1)=0\) and \(r_2(1)=1\). At \(N=4\), (10) and
(15) give the profile \((1,1,1)\). Every nontrivial Fourier matrix has a
nonzero corner at \(k=1\), so the minima at \(N=2,4\) are exactly one.
This also explains why substituting \(N=4\) into \(N/4+1\) is invalid.

\subsection{2.4 Half squares and the characteristic zero
extension}\label{half-squares-and-the-characteristic-zero-extension}

For every complex Fourier rectangle \(V=U_N[R,H]\), coordinate
restriction gives \(\|V\|_2\leq1\), while \(\|V\|_F^2=|R||H|/N\). The
sum of the squared singular values is at most their number, hence

\[
\operatorname{rank}V\geq |R||H|/N.                         \tag{16}
\]

For \(N=4L\) and half-size sets, this is \(\operatorname{rank}V\geq L\).
Choose both sets to be all even labels. With \(0\leq a,b<2L\), the
unnormalized entries are

\[
\omega_N^{(2a)(2b)}=\omega_L^{ab}.
\]

Every residue modulo \(L\) occurs twice in each axis, so the block is
\(F_L\) with repeated rows and columns and has rank exactly \(L\). This
proves (5). The chosen set has a nontrivial translation stabilizer,
consistent with deficit zero in Theorem 2. At \(N=2\), every half-square
is a nonzero scalar.

Finally let \(K\) have characteristic zero and contain a primitive
\(N\)th root \(\zeta\). Every matrix involved in the exact statements
has entries in the cyclotomic subfield \(\mathbb Q(\zeta)\). Matrix rank
is unchanged under field extension, since it is characterized by
vanishing and nonvanishing minors. An embedding of this subfield into
\(\mathbb C\) sends \(\zeta\) to a primitive complex root. Such a root
is \(\omega_N^a\) with \(\gcd(a,N)=1\), giving the same minimax through
a row-label permutation. The bit-reversal recurrence also holds directly
for every primitive root. This proves the asserted field extension
without assuming that all of \(K\) embeds into \(\mathbb C\) or that
\(\sqrt N\in K\). It does not transfer the positivity proof to finite
characteristic. This completes Theorem 1.

\section{3 The finite abelian autocorrelation deficit
theorem}\label{the-finite-abelian-autocorrelation-deficit-theorem}

\textbf{Theorem 2.} Let \(G\) be any finite abelian group of order
\(N\geq2\), and let \(U\) be (3). Fix \(R\subseteq\widehat G\) and
\(H\subseteq G\) with \(1\leq r=|R|,h=|H|\leq N\). Put \(V=U[R,H]\) and

\[
\begin{aligned}
c_H(s)&=|H\cap(H-s)|,\\
d_H&=h-\max_{s\neq0}c_H(s),\\
B_H&=\frac{h(N-r)}N-\frac{d_Hr(N-r)}{N^2}.
\end{aligned}                                                \tag{17}
\]

Then

\[
\operatorname{rank}V\geq
\left\lceil\frac{rh}N+\frac{d_Hr(N-r)}{N^2}\right\rceil.
                                                               \tag{18}
\]

For every integer \(0\leq t<h-B_H\) and every
\(W\in\mathbb C^{r\times h}\) with \(\operatorname{rank}W\leq t\),

\[
\|V-W\|_2\geq\sigma_{t+1}(V)
\geq\left(1-\sqrt{\frac{B_H}{h-t}}\right)^{1/2}>0.             \tag{19}
\]

The transposed bounds hold with the roles of \(R\) and \(H\)
interchanged. Explicitly, define \(d_R\) by nonzero translations in
\(\widehat G\) and put

\[
B_R=\frac{r(N-h)}N-\frac{d_Rh(N-h)}{N^2}.
\]

Then the other rank bound is

\[
\operatorname{rank}V\geq
\left\lceil\frac{rh}N+\frac{d_Rh(N-h)}{N^2}\right\rceil,       \tag{20}
\]

and the other gap is \((1-\sqrt{B_R/(r-t)})^{1/2}\) when
\(0\leq t<r-B_R\). One may take the larger rank lower bound and, among
gap bounds whose own strict condition holds, the larger gap.

\subsection{3.1 Proof of the rank
estimate}\label{proof-of-the-rank-estimate}

Let \(I=\widehat G\setminus R\) and \(n=N-r\). Set

\[
A=U[I,H]^*U[I,H],\qquad f_I(s)=\sum_{\chi\in I}\chi(s).
\]

As in (6), \(A\) is a positive contraction and \(V^*V=I_h-A\). For
\(x,y\in H\), \(A_{xy}=N^{-1}f_I(y-x)\). Thus the exact second-moment
identity and orthogonality give

\[
\operatorname{tr}(A^2)=\frac1{N^2}\sum_{s\in G}c_H(s)|f_I(s)|^2,
\quad \sum_{s\in G}|f_I(s)|^2=Nn,\quad f_I(0)=n.              \tag{21}
\]

The identity uses \(\operatorname{tr}(A^2)=\sum_{x,y}|A_{xy}|^2\), valid
because \(A\) is Hermitian. Using \(c_H(0)=h\) and \(c_H(s)\leq h-d_H\)
for \(s\neq0\) yields

\[
\begin{aligned}
\operatorname{tr}(A^2)
&\leq\frac{h n^2+(h-d_H)(Nn-n^2)}{N^2}\\
&=\frac{hn}N-\frac{d_Hrn}{N^2}=B_H.                         \end{aligned} \tag{22}
\]

The multiplicity of eigenvalue one of \(A\) equals \(\dim\ker V\) and is
at most \(\operatorname{tr}(A^2)\). Hence
\(\operatorname{rank}V\geq h-B_H\). Taking the integer ceiling proves
(18). Also \(0\leq d_H\leq h\), so \(B_H=(n/N)(h-d_Hr/N)\geq0\), as
required for (19).

\subsection{3.2 Proof of the indexed gap and the transpose
statement}\label{proof-of-the-indexed-gap-and-the-transpose-statement}

List the eigenvalues of \(A\) as \(a_1\geq\cdots\geq a_h\), all in
\([0,1]\). The descending eigenvalues of \(V^*V\) are
\(1-a_h,\ldots,1-a_1\). Therefore

\[
\sigma_{t+1}(V)^2=1-a_{h-t}.
\]

Since at least \(h-t\) eigenvalues of \(A\) are at least \(a_{h-t}\),

\[
(h-t)a_{h-t}^2\leq\sum_{j=1}^h a_j^2
=\operatorname{tr}(A^2)\leq B_H.
\]

This proves the second inequality in (19). The strict condition
\(t<h-B_H\) implies \(h-t>B_H\geq0\), so the denominator is positive and
the resulting gap is positive. This condition and (18) also ensure
\(t<\operatorname{rank}V\leq\min(r,h)\), so the stated nonzero
singular-value index is legitimate.

For completeness, let \(E\) be the span of the top \(t+1\) right
singular vectors of \(V\). Since \(\operatorname{rank}(W|_E)\leq t\),
some unit vector \(x\in E\) satisfies \(Wx=0\). Then
\(\|(V-W)x\|_2=\|Vx\|_2\geq\sigma_{t+1}(V)\), proving the first
inequality in (19) without a separate approximation theorem as an
unproved premise.

Under the canonical identification \(G\cong\widehat{\widehat G}\), the
transpose \(U^T\) is the unitary character table of \(\widehat G\).
Apply the proved statements to \(V^T\), whose rank, nonzero singular
values, and operator norm equal those of \(V\). This proves (20) and the
transposed gap. Theorem 2 is complete.

\subsection{3.3 Endpoint and normalization
cases}\label{endpoint-and-normalization-cases}

The definition of \(d_H\) does not apply to the trivial group, since
there is no nonzero translation. At \(N=1\) the character table is
\([1]\): its unique nonempty block has rank one and distance one from
the only rank-zero matrix. Empty restrictions have rank zero. These
direct facts are the entire \(N=1\) convention, not an artificial
maximum over an empty set.

For \(N\geq2\), empty \(R\) or \(H\) may be handled directly as
rank-zero blocks; they are excluded from the statement to avoid
unnecessary deficit and approximation-index conventions. When \(r=N\),
\(B_H=0\) and \(V\) has orthonormal columns. Equations (18) and (19)
give rank \(h\) and gap one for every \(0\leq t<h\). When \(h=N\),
\(d_H=0\) and (18) gives rank \(r\); the transposed bound gives the
exact gap one for every \(t<r\), since \(V\) has orthonormal rows.

For a singleton \(H\), \(d_H=1\) and \(B_H=(N-r)^2/N^2\). The rank bound
rounds up to one and the gap at \(t=0\) is exactly \(\sqrt{r/N}\), the
norm of this column. For a singleton \(R\), the transposed statement
gives the exact norm \(\sqrt{h/N}\). These observations cover all
nonempty restrictions at \(N=2\) as well as singleton cases at larger
sizes.

If \(t\geq h-B_H\), equation (19) asserts no strictly positive gap from
that orientation. Equality at the threshold gives zero in the displayed
expression; when \(t\geq h\) the denominator is not used. Failure of
this sufficient strict condition is not a claim that an exact rank-\(t\)
representation exists. The other orientation may still apply.

For the unnormalized character table \(F=\sqrt N\,U\), all rank
conclusions are unchanged. For a rank-at-most-\(t\) approximant \(Z\) to
\(F[R,H]\), apply (19) to \(W=Z/\sqrt N\): the lower bound on
\(\|F[R,H]-Z\|_2\) is \(\sqrt N\) times the normalized bound. More
generally, multiplication of the whole target by a scalar \(\alpha\)
multiplies the norm lower bound by \(|\alpha|\). This is the
normalization needed for every gap in this manuscript.

\section{4 Existing consequences and their
boundaries}\label{existing-consequences-and-their-boundaries}

\subsection{Translation stabilizers and
cardinality}\label{translation-stabilizers-and-cardinality}

For \(s\neq0\), \(c_H(s)=h\) if and only if \(H+s=H\). Thus \(d_H>0\)
exactly when \(H\) has trivial translation stabilizer. If \(G\) is a
\(p\)-group and \(p\nmid h\), each nonidentity translation has orbit
length a positive power of \(p\), so \(H\) cannot be a union of its
orbits. It follows that \(d_H\geq1\). More generally, \(\gcd(h,N)=1\)
suffices: the nontrivial order of a translation divides \(N\) and would
have to divide \(h\) for invariance.

For a square block of side \(s\) in a \(p\)-group with \(p\nmid s\),
(18) consequently gives

\[
\operatorname{rank}U[R,H]\geq
\left\lceil\frac{s^2}N+\frac{s(N-s)}{N^2}\right\rceil.        \tag{23}
\]

For example, at \(N=64\) and \(s=21\), the expression inside the ceiling
is \(29127/4096\), giving eight; the ordinary ceiling
\(\lceil441/64\rceil\) is seven. No location assumption or new separator
model is needed for this existing example.

\subsection{Critical rectangles for abelian two
groups}\label{critical-rectangles-for-abelian-two-groups}

For every finite abelian \(2\)-group of order \(N=4L\geq8\), take
\(r=2L-1\), \(h=2L+1\). Odd cardinality gives \(d_H\geq1\) and (22)
gives the same upper bound \(B_L\) in (8). With \(t=L\), (19) yields the
existing corollary

\[
\sigma_{L+1}(U[R,H])\geq
\delta_N:=\left(1-\sqrt{\frac{B_L}{L+1}}\right)^{1/2}>0.      \tag{24}
\]

Indeed, \(L<h-B_L\) is equivalent to \(B_L<L+1\). Since

\[
L+1-B_L=\frac14-\frac1{4L}-\frac1{16L^2}\longrightarrow\frac14,
\]

the expansion \(1-\sqrt{1-x}\sim x/2\) gives
\(\delta_N\sim(2N)^{-1/2}\). This is a decaying indexed gap, not uniform
conditioning of the entire rectangle. If a full-matrix approximant has
rank at most \(L\) on the corresponding critical corner, its
operator-norm error from the independently reordered normalized
character table is at least \(\delta_N\), since coordinate restriction
cannot increase the error norm. A matching endpoint minimax upper
construction for every noncyclic \(2\)-group is not supplied.

\subsection{Rank sums at the already accepted
sizes}\label{rank-sums-at-the-already-accepted-sizes}

Let \(S(M)=\sum_{k=1}^{N-1}\operatorname{rank}M[0:k,k:N]\). Equation
(16) gives the universal pointwise lower bound
\(\lceil k(N-k)/N\rceil\). For odd \(k\) and dyadic \(N\geq2\), the
column size \(h=N-k\) is odd, so (18) improves this to

\[
\operatorname{rank}M[0:k,k:N]\geq
\left\lceil\frac{k(N-k)(N+1)}{N^2}\right\rceil.              \tag{25}
\]

At \(N=8\), these lower bounds are \((1,2,3,2,3,2,1)\). At \(N=16\),
they are \((1,2,3,3,4,4,5,4,5,4,4,3,3,2,1)\). Recurrence (10) attains
each entry. The exact minimum upper-cut sums at these sizes are
therefore 14 and 48. The base-size minima are \(S_1=0\), \(S_2=1\), and
\(S_4=3\).

For the bit-reversal ordering only, summing (10) using reflection gives

\[
S_N=\frac{N^2}{8}+\frac N4+4S_{N/4},\qquad S_1=0, S_2=1.
\]

To see the coefficients, the even cuts up to the middle contribute
\(\sum_{a=1}^{L}(a+r_L(a))\), the odd cuts contribute
\(\sum_{a=0}^{L-1}(a+1+r_L(a))\), and reflection doubles the combined
sum and subtracts the middle value \(r_N(2L)=L\). For \(N=2^m\), the
solution is

\[
S_N=\begin{cases}
N^2/6+N(m/8-1/6),&m\text{ even},\\
N^2/6+N(m/8+1/24),&m\text{ odd}.
\end{cases}                                                  \tag{26}
\]

Meanwhile summing the unrounded bound in (16) gives
\(S(M)\geq(N^2-1)/6\). Thus (26) has the optimal leading term \(N^2/6\).
It does not prove exact all-size minimum rank sums. Already at \(N=32\),
(25) together with the even-cut baseline sums to 182, whereas (26) gives
192; these bounds do not settle the gap.

\section{5 Prior art and the recovered final paper
comparison}\label{prior-art-and-the-recovered-final-paper-comparison}

Double digit reversal and the simultaneous rank-one block hierarchy are
established in Delvaux and Van Barel's final 2008 article {[}1{]},
especially Eq. (15) and Lemma 7. The upper calculation in Section 2 is
explicit cut bookkeeping within that known structure. Their Theorem 9
describes the minimum Fourier support of one nonzero vector with bounded
input support; Remark 10 gives additional extremal-support geometry. The
\(N=8\) critical lower bound is already a consequence: a rank-deficient
\(5\times3\) transposed critical block would give input support at most
three and output support at most three, whereas their theorem requires
at least four.

For \(N=4L\geq16\), the direct one-vector substitution at a hypothetical
rank \(L\) only gives Fourier support at least four, while the two
complement sizes are \(2L-1\) and \(2L+1\). This direct substitution
does not exclude rank \(L\). That observation distinguishes quantifiers;
it is not a proof that the entire older theory cannot imply the general
result through a different argument.

The final comparison report {[}S6{]} has now been recovered
byte-for-byte at its recorded hash, and the local final journal PDF
matches the hash recorded in that closure. The relevant final passages
were checked again against this PDF. Its source-availability issue is
closed. Older immutable reports {[}S2, S5{]} contain superseded access
wording; that wording is not a current blocker. The bounded comparison
concluded that the complete general minimax and deficit rank/gap
statements were not stated in that article. It did not establish
worldwide novelty.

The ordinary flat-unitary bound (16) is classical; de Wolf {[}4, Theorem
1{]} gives the Hadamard version, attributes it to Lokam, and discusses
the equal-magnitude orthogonal/DFT setting. Projection and
Hilbert--Schmidt uncertainty methods are classical {[}7{]}.
Stabilizer-sensitive support uncertainty is also established {[}5, Lemma
7.1 and Theorem 7.2{]}. The statistic used here is not newly invented:
Shkredov {[}6, Eq. (4){]} uses \(\rho(H)=\max_{s\neq0}|H\cap(H+s)|\),
and exactly \(d_H=h-\rho(H)\). The sign of the translation has no effect
on this maximum.

The Kronecker Fourier work {[}3{]} is relevant prior art for
subgroup-based rank-deficient blocks and support extrema on general
finite abelian groups. Merely extending the group class is not an
originality argument. The remaining historical comparison concerns the
exact overlap-retaining fixed-rectangle coefficient, higher-dimensional
support or rank profiles, and concentration second moments. Those
questions are separate from the complete proofs above.

\subsection{Versioned citation
concordance}\label{versioned-citation-concordance}

\begin{longtable}[]{@{}
  >{\raggedright\arraybackslash}p{(\columnwidth - 4\tabcolsep) * \real{0.2600}}
  >{\raggedright\arraybackslash}p{(\columnwidth - 4\tabcolsep) * \real{0.3800}}
  >{\raggedright\arraybackslash}p{(\columnwidth - 4\tabcolsep) * \real{0.3600}}@{}}
\toprule\noalign{}
\begin{minipage}[b]{\linewidth}\raggedright
Source and version
\end{minipage} & \begin{minipage}[b]{\linewidth}\raggedright
Hierarchy or support passages
\end{minipage} & \begin{minipage}[b]{\linewidth}\raggedright
Verification in this packet
\end{minipage} \\
\midrule\noalign{}
\endhead
\bottomrule\noalign{}
\endlastfoot
DV2008 final {[}1{]} & Eq. (15), pp.~1593--1594; Lemma 7, p.~1596;
Theorem 9, pp.~1598--1599; Remark 10, p.~1599 & Recovered accepted
comparison; matching local final PDF; key passages rechecked \\
TW470 September 2006 {[}2{]} & Lemma 8, printed p.~11 / PDF p.~13;
Theorem 10, printed pp.~13--14 / PDF pp.~15--16 & Author-institution PDF
and numbering freshly checked \\
TW477 November 2006 {[}3{]} & Theorem 14, printed pp.~11--12; Theorem
17, p.~13; Theorem 20 and Corollary 21, pp.~17--18 & Author-institution
preprint; earlier bounded comparison retained; no relabeling as
final-journal numbering \\
\end{longtable}

The preprint and final article numbers are deliberately not
interchangeable. The local inspection did not repeat the entire
final-paper comparison or all earlier specialist literature searches.
References {[}4{]}--{[}7{]} were checked at their indicated
primary-source versions; the broader inherited reference inventory is
preserved in the source packet.

\section{6 Status and open extensions}\label{status-and-open-extensions}

No mathematical gap was found in consolidating the two stated theorems.
Their previously recorded independent checks are {[}S3{]} and {[}S4{]}.
The present manuscript is an author-side consolidation awaiting a separate independent manuscript review; it does
not label itself newly independently accepted. The algebraic controls
accompanying the packet are regression checks, not replacements for the
proofs or new finite-field theorems.

The following are optional future problems, not omitted assumptions or
unfinished steps in either proof:

\begin{itemize}
\tightlist
\item
  Find matching endpoint minimax upper constructions for arbitrary
  noncyclic finite abelian \(2\)-groups.
\item
  Determine the exact minimum upper-cut rank sum at every dyadic size
  beyond the accepted small cases.
\item
  Determine sharper or optimal fixed-rectangle gap bounds, or suitable
  distinct positive-characteristic rank statements.
\item
  Complete the wider model-matched historical comparison for the general
  minimax and the exact deficit formula.
\end{itemize}

No separator fusion is developed in this manuscript. Nor does high cut
rank alone measure Fourier evaluation cost: even a permutation can have
large corner rank while applying cheaply. Arithmetic or circuit lower
bounds would require a separately specified and charged computation
model. Neither theorem supplies one.

\section{References}\label{references}

\begin{enumerate}
\def\labelenumi{\arabic{enumi}.}
\tightlist
\item
  Steven Delvaux and Marc Van Barel. \emph{Rank-deficient submatrices of
  Fourier matrices}. Linear Algebra and its Applications 429 (2008),
  1587--1605.
  \href{https://doi.org/10.1016/j.laa.2008.04.043}{doi:10.1016/j.laa.2008.04.043}.
  Final journal version inspected locally; source PDF SHA-256 is
  recorded in the packet.
\item
  Steven Delvaux and Marc Van Barel. \emph{Rank-deficient submatrices of
  Fourier matrices}. Report TW470, September 2006.
  \href{https://www.cs.kuleuven.be/publicaties/rapporten/tw/TW470.pdf}{Author-institution
  preprint}. Numbering is specific to this version.
\item
  Steven Delvaux and Marc Van Barel. \emph{Rank-deficient submatrices of
  Kronecker products of Fourier matrices}. Report TW477, November 2006.
  \href{https://www.cs.kuleuven.be/publicaties/rapporten/tw/TW477.pdf}{Author-institution
  preprint}. The later article is in Linear Algebra and its Applications
  426 (2007), 349--367; passage references here use TW477.
\item
  Ronald de Wolf. \emph{Lower Bounds on Matrix Rigidity via a Quantum
  Argument}.
  \href{https://homepages.cwi.nl/~rdewolf/publ/qc/qrigid.pdf}{Author-hosted
  preprint}, Theorem 1, p.~4, and Section 1.2, p.~2. The flat-matrix
  lower bound is credited there to Lokam; no unverified original
  corollary number is asserted here.
\item
  Tao Feng, Henk D. L. Hollmann, and Qing Xiang. \emph{The shift bound
  for abelian codes and generalizations of the Donoho--Stark uncertainty
  principle}.
  \href{https://arxiv.org/pdf/1804.00367v1}{arXiv:1804.00367v1}, 2018,
  Lemma 7.1 and Theorem 7.2, PDF p.~11.
\item
  Ilya D. Shkredov. \emph{Uncertainty for convolutions of sets}.
  \href{https://arxiv.org/pdf/2404.12469v1}{arXiv:2404.12469v1}, 2024,
  Eq. (4), p.~2. The cited definition is the maximum nonzero-translation
  overlap.
\item
  Gorjan Alagic and Alexander Russell. \emph{Uncertainty Principles for
  Compact Groups}.
  \href{https://arxiv.org/pdf/math/0608702v2}{arXiv:math/0608702v2},
  2008, Theorem 1 and Corollary 1, pp.~1--2, and Theorem 3.1.1, p.~4.
\end{enumerate}

\section{Appendix Proof source
concordance}\label{appendix-proof-source-concordance}

The following labels refer to exact local source copies in the
accompanying proof-source packet. Complete byte counts, original paths,
retrieval provenance, and full SHA-256 hashes are in
SOURCE\_MANIFEST.json. The source copies remain unchanged, including
their historical dating and now-superseded access prose.

\begin{longtable}[]{@{}
  >{\raggedright\arraybackslash}p{(\columnwidth - 4\tabcolsep) * \real{0.0900}}
  >{\raggedright\arraybackslash}p{(\columnwidth - 4\tabcolsep) * \real{0.4100}}
  >{\raggedright\arraybackslash}p{(\columnwidth - 4\tabcolsep) * \real{0.5000}}@{}}
\toprule\noalign{}
\begin{minipage}[b]{\linewidth}\raggedright
Label
\end{minipage} & \begin{minipage}[b]{\linewidth}\raggedright
Source
\end{minipage} & \begin{minipage}[b]{\linewidth}\raggedright
Role in this manuscript
\end{minipage} \\
\midrule\noalign{}
\endhead
\bottomrule\noalign{}
\endlastfoot
S1 & MINIMAX\_PROOF, \emph{Independent Fourier endpoint rank audit} &
Theorem 1, exact recurrence, small rank sums \\
S2 & DEFICIT\_PROOF, \emph{Autocorrelation-deficit Fourier bounds and
routed separators} & Theorem 2, stabilizer/cardinality consequences;
separator sections excluded here \\
S3 & R07, REVIEW-07 report, Section 7 C02 & Recorded independent
acceptance of endpoint scope and critical rectangle extension \\
S4 & R08, REVIEW-08 report, deficit passage & Recorded independent
checking of the separate deficit estimate and spectral index \\
S5 & MINIMAX\_PRIORITY, endpoint priority report & Half-square boundary,
prior-art boundaries, version-specific references \\
S6 & Recovered DELVAUX\_FINAL\_COMPARISON & Half-square proof and
completed bounded final-paper comparison \\
S7 & FINAL\_THEOREMS, Section 1 & Existing abelian two-group
critical-rectangle gap and detailed upper proof \\
S8 & ORIGINAL\_MINIMAX\_PROOF & Original exact recurrence exposition
extracted inertly from its uploaded archive \\
\end{longtable}

The five readiness-pinned main sources were rehashed and matched before
writing. Source S6 was recovered from the repository blob already listed
in the local pinned tree; it matches its inherited comparison hash
exactly. The packet does not redistribute the journal article PDF or its
extracted full text. It records the verified local PDF hash and its
bibliographic reference instead.

\end{document}
